%-----------------------------------------------------------------------------------------------
\makeatletter
\immediate\write18{datelog > \jobname.info} % site script for $(date '+%Y-%m-%d %Hh%Mm%Ss')
\makeatother
%-----------------------------------------------------------------------------------------------
%-----------------------------------------------------------------------------------------------
\usetheme{Copenhagen}
\usepackage{beamercolorthemeCNThermSci2}
\usefonttheme{serif}
%-----------------------------------------------------------------------------------------------

%-----------------------------------------------------------------------------------------------
%-----------------------------------------------------------------------------------------------
\usetheme{Copenhagen}
\usepackage{beamercolorthemeUTF2}
\usefonttheme{serif}
%-----------------------------------------------------------------------------------------------
\usepackage[utf8]{inputenc}
\usepackage[greek,french,english,brazil]{babel} % last becomes the active one
\usepackage{pslatex}
\usepackage{amssymb,amsmath}
\usepackage{soul}
\usepackage[squaren,Gray,cdot]{SIunits}
\usepackage[nice]{nicefrac}
\usepackage{tikz}
\usepackage{amscd}
\usepackage{stmaryrd}
\usepackage{scalerel}
\usepackage{xspace}
%-----------------------------------------------------------------------------------------------


%-----------------------------------------------------------------------------------------------
%-----------------------------------------------------------------------------------------------
% Mathematical
%-----------------------------------------------------------------------------------------------
\newcommand{\vet}[1]{\underline{{#1}}}
\newcommand{\mat}[1]{\underline{\underline{{#1}}}}
\newcommand{\cub}[1]{\underline{\underline{\underline{{#1}}}}}
\newcommand{\eqdef}{{\ensuremath\stackrel{\text{\tiny def}}{=}}}
%-----------------------------------------------------------------------------------------------
% Linguistic
%-----------------------------------------------------------------------------------------------
\newcommand{\GRtxt}[1]{\begin{otherlanguage}{greek}{{#1}}\end{otherlanguage}}
\newcommand{\FRtxt}[1]{\begin{otherlanguage}{french}{{#1}}\end{otherlanguage}}
%-----------------------------------------------------------------------------------------------
% Presentation
%-----------------------------------------------------------------------------------------------
\newcommand{\BkgImgH}[1]{% Places an image centered on the slide background filling the height
    \usebackgroundtemplate{\parbox{\paperwidth}{%
        \vspace*{1sp}\centering\includegraphics[height=\paperheight]{{#1}}
}}}
\newcommand{\BkgImgW}[1]{% Places an image centered on the slide background filling the width
    \usebackgroundtemplate{\parbox{\paperwidth}{%
        \vspace*{1sp}\centering\includegraphics[width=\paperwidth]{{#1}}
}}}
\newcommand{\ArtEndH}[3]{% Transitions to plain image (last) slide: #1:prefix #2,#3:extensions
    \BkgImgH{root/../art/#1.#2}
    \frame<handout:0>[plain]{%
        \transdissolve\vspace*{72mm}\color{white}\scriptsize\bf\input{root/../art/#1.#3}}
    \usebackgroundtemplate{\mbox{~}}
}
\newcommand{\ArtEndW}[3]{% Transitions to plain image (last) slide: #1:prefix #2,#3:extensions
    \BkgImgW{root/../art/#1.#2}
    \frame<handout:0>[plain]{%
        \transdissolve\vspace*{72mm}\color{white}\scriptsize\bf\input{root/../art/#1.#3}}
    \usebackgroundtemplate{\mbox{~}}
}
\newcommand{\ImgColW}[3]{% Inserts a full-width image in a column
    \includegraphics[width=\columnwidth]{root/../art/#1.#2}\\[-0.5\baselineskip]
    \parbox{\columnwidth}{\tiny\hfill\scalebox{0.85}{\input{root/../art/#1.#3}}}
}
\newcommand{\txtpic}[1]{%
    \fcolorbox{lightgray}{white!90!black}{{#1}} 
}
%-----------------------------------------------------------------------------------------------


%-----------------------------------------------------------------------------------------------
\title{C.02.01.A2 -- Combustão e Equilíbrio Químico}
\subtitle{Aplicação em FTAF -- Finite Time Air-Fuel Otto Engine Model}
\author{Prof.~C.~Naaktgeboren, PhD}
\date{{\scriptsize\tt%
    \includegraphics[height=6.0mm]{root/00-res/cc/by-nc-nd-88x31.pdf}\\[\smallskipamount]
    https://github.com/CNThermSci/ApplThermSci\\
    Compiled on \input{\jobname.info}
}}
%-----------------------------------------------------------------------------------------------
\begin{document}
%-----------------------------------------------------------------------------------------------
\logo{%
    \parbox{158mm}{% There's a 1mm gap on each side of the 160mm x 90mm slide logo line
        \mode<beamer>{
            \includegraphics[height=6.0mm]{root/00-res/UTFPR/UTFPR-logo-D.pdf}\hfill%
            \includegraphics[height=9.0mm]{root/00-res/logo/CNThermSci-logo-A.pdf}%
        }
        \mode<handout>{
            \includegraphics[height=6.0mm]{root/00-res/UTFPR/UTFPR-logo-W.pdf}\hfill%
            \includegraphics[height=9.0mm]{root/00-res/logo/CNThermSci-logo-W.pdf}%
        }
    }
} % The (delineated, alpha), or washed-out logos
%-----------------------------------------------------------------------------------------------
\frame{\titlepage}
%-----------------------------------------------------------------------------------------------

%----------------------------r-------------------------------------------------------------------
\section{Modelos de Misturas Reagentes}
%-----------------------------------------------------------------------------------------------

%----------------------------r-------------------------------------------------------------------
\subsection{Modelo de Ar}
%-----------------------------------------------------------------------------------------------

    % !j 96 -i8
    %-------------------------------------------------------------------------------------------
    \begin{frame}{Modelo de Ar}\vspace*{-2em}
        %-----
        \begin{itemize}
            \item<1-> Ar é modelado \alert{apenas} como uma mistura de Oxigênio,
                \alert{\ce{O2}}, e Nitrogênio, \alert{\ce{N2}}; \\[\bigskipamount]
            \item<2-> A proporção é de \alert{$\unit{\psi}{\kilo\mole}$ de \ce{N2}} para cada
                \alert{$\unit{1}{\kilo\mole}$ de \ce{O2}}; \\[\bigskipamount]
            \item<3-> Nitrogênio será considerado \alert{gás inerte}; \\[\bigskipamount]
            \item<4-> Todos os demais gases inertes são modelados como sendo \alert{\ce{N2}};
                \\[\bigskipamount]
            \item<5-> Valor típico para \alert{$\psi$} é de \alert{$79/21 \approx 3,76$}.
        \end{itemize}
    \end{frame}
    %-------------------------------------------------------------------------------------------

%-----------------------------------------------------------------------------------------------
\subsection{Modelo de Vapor de Combustível}
%-----------------------------------------------------------------------------------------------

    % !j 96 -i8
    %-------------------------------------------------------------------------------------------
    \begin{frame}{Modelo de Combustível}\vspace*{-2em}
        %-----
        \begin{itemize}
            \item<1-> Molécula de combustível \alert{modelada} como \alert{\ce{C\NC H\NH O\NO
                N\NN}}; \\[\bigskipamount]
            \item<2-> Valores \alert{\NC}, \alert{\NH}, \alert{\NO}, e \alert{\NN} são
                \alert{parâmetros} ajustáveis; \\[\bigskipamount]
            \item<3-> Seja \alert{$\epsilon$} a quantidade de \alert{combustível por kmol de
                \ce{O2}} estequiometricamente oxidada;
        \end{itemize}
        \vspace*{0.8em}\begin{equation*}
            \uncover<4->{
                \alert{\epsilon^{-1} \equiv \NC + \frac{\NH}{4} - \frac{\NO}{2}}.
            }
        \end{equation*}
        %-----
        \vspace*{-0.8em}\begin{itemize}
            \item<5-> \alert{$\epsilon/(1+\psi)$} é a razão \alert{combustível-ar}
                estequiométrica.
        \end{itemize}
    \end{frame}
    %-------------------------------------------------------------------------------------------

    % !j 96 -i8
    %-------------------------------------------------------------------------------------------
    \begin{frame}{Razão de Equivalência:}\vspace*{-2em}
        %-----
        \begin{itemize}
            \item<1-> Seja \alert{$\phi$} a razão de equivalência, ou a razão combustível-ar
                \alert{real pela estequiométrica}:
        \end{itemize}
        \vspace*{0.8em}\begin{equation*}
            \uncover<2->{
                \alert{\phi \equiv \frac{\NF/\NA}{\epsilon / (1 + \psi)}},\qquad\mbox{assim,}
            }
        \end{equation*}
        %-----
        \vspace*{-0.8em}\begin{itemize}
            \item<3-> \alert{$\phi < 1$} modela misturas combustível-ar com
                \alert{excesso de ar} (pobre em combustível);
                \\[\bigskipamount]
            \item<4-> \alert{$\phi > 1$} modela misturas combustível-ar com
                \alert{excesso de combustível} (pobre em ar); e
                \\[\bigskipamount]
            \item<5-> \alert{$\phi = 1$} modela misturas combustível-ar
                \alert{estequiométricas}.
        \end{itemize}
    \end{frame}
    %-------------------------------------------------------------------------------------------

    % Finishes with stunning image, with credit
    \ArtEndWSE{pexels-eberhard-grossgasteiger-542608}{jpg}{txt}

%-----------------------------------------------------------------------------------------------
\end{document}
%-----------------------------------------------------------------------------------------------

